Coefficient priors can affect the prediction, interpretation and computational reliability of Bayesian survival models. This thesis compares three prior specifications with Cox proportional-hazards regression, random survival forests and DeepSurv as reference methods. It brings held-out prediction, calibration, posterior uncertainty and explanation stability into a controlled evaluation of clinical survival modelling.

Three cohorts were analysed: GBSG2 (686 patients), METABRIC (1,980) and WHAS500 (500). The Bayesian models shared a piecewise-exponential proportional-hazards likelihood with six baseline intervals and differed in their Normal, regularised Horseshoe or continuous Spike-and-Slab coefficient priors. Evaluation used five repeats of five-fold cross-validation, training-fold preprocessing, inner tuning for random survival forests and DeepSurv, and predefined Bayesian diagnostic criteria.

Of 450 model--fold outcomes, 430 met the evaluation criteria, including 205 of 225 Bayesian fits. Mean integrated Brier scores for Normal, continuous Spike-and-Slab and Cox regression were 0.1672, 0.1672 and 0.1673 in GBSG2. Random survival forests had the lowest mean error among models with complete fold coverage in GBSG2 (0.1624) and METABRIC (0.1806). Bayesian scores ranged from 0.1521 to 0.1524 in WHAS500. Five of nine primary prior contrasts had complete accepted pairs; their corrected 95\% confidence intervals all contained zero. Matched-fold posterior summaries showed close Normal and continuous-mixture estimates, with more selective Horseshoe shrinkage. Calibration identified early-risk underprediction by the Bayesian models in WHAS500.

Robustness analyses accepted 770 of 825 model--condition--fold outcomes. Posterior contribution and exact linear SHAP rankings agreed closely, while 102 of 1,640 stability comparisons fell outside tolerance. A separate post-hoc sampler sensitivity accepted all 75 Horseshoe fits and detected no predictive advantage between priors. The findings show how prior choice can change coefficient interpretation and sampling reliability even when predictive differences are small, providing practical evidence for evaluating Bayesian survival models on complementary criteria.

Keywords: Bayesian survival analysis; proportional hazards; shrinkage priors; calibration; explanation stability; repeated cross-validation.
